% Packages for a PDFLaTeX-compatible presentation on Overleaf.
\usepackage{cmap}
\usepackage[utf8]{inputenc}
\usepackage[T1]{fontenc}
\usepackage[main=english]{babel}
\babelprovide[import]{brazilian}
\babelprovide[import]{spanish}
\usepackage{xpatch}
\usepackage{pgfpages}
\usepackage{fancyvrb}
\usepackage{tikz}
\usepackage{booktabs}
\usepackage{colortbl}
\usepackage{tabularx}
\usepackage{subcaption}
\usepackage[authoryear,round]{natbib}

% Typography and slide margins.
\usefonttheme{serif}
\setbeamersize{text margin left=14mm,text margin right=14mm}

\setbeamerfont{title}{size=\Huge,series=\bfseries}
\setbeamerfont{subtitle}{size=\large,series=\mdseries}
\setbeamerfont{author}{size=\small,series=\mdseries}
\setbeamerfont{institute}{size=\footnotesize}
\setbeamerfont{frametitle}{size=\large,series=\bfseries}
\setbeamerfont{framesubtitle}{size=\small,series=\mdseries}
\setbeamerfont{part title}{size=\Large,series=\mdseries}
\setbeamerfont{caption}{size=\footnotesize}
\setbeamerfont{caption name}{series=\bfseries}
\setbeamerfont{footline}{size=\fontsize{6}{7}\selectfont}
\setbeamerfont{footnote}{size=\tiny}
\setbeamerfont{block title}{size=\normalsize,series=\bfseries}
\setbeamerfont{block body}{size=\normalsize}
\setbeamerfont{itemize/enumerate subbody}{size=\normalsize}
\setbeamerfont{section in toc}{size=\normalsize,series=\bfseries}
\setbeamerfont{subsection in toc}{size=\small}
\setbeamerfont{subsubsection in toc}{size=\footnotesize}

% Shared palette: off-white, soft black, and blue accents.
\definecolor{offwhite}{HTML}{F8F8F8}
\definecolor{softblack}{HTML}{282828}
\definecolor{darkblue}{HTML}{324B64}
\definecolor{lightblue}{HTML}{9EC4D9}
\definecolor{vibrantblue}{HTML}{327DE1}
\newcommand{\keyphrase}[1]{\textcolor{lightblue}{#1}}
\newcommand{\highlightrow}{\rowcolor{darkblue}}
\tikzset{
    progress arrow/.style={->,draw=vibrantblue},
    final stage/.style={draw=vibrantblue}
}

\setbeamercolor{normal text}{fg=offwhite,bg=softblack}
\setbeamercolor{background canvas}{bg=softblack}
\setbeamercolor{palette primary}{fg=offwhite,bg=softblack}
\setbeamercolor{palette secondary}{fg=offwhite,bg=softblack}
\setbeamercolor{palette tertiary}{fg=offwhite,bg=softblack}
\setbeamercolor{palette quaternary}{fg=offwhite,bg=softblack}
\setbeamercolor{structure}{fg=offwhite}
\setbeamercolor{item}{fg=offwhite}
\setbeamercolor{title}{fg=offwhite}
\setbeamercolor{author}{fg=offwhite}
\setbeamercolor{institute}{fg=offwhite}
\setbeamercolor{frametitle}{fg=offwhite}
\setbeamercolor{framesubtitle}{fg=offwhite}
\setbeamercolor{page number in head/foot}{fg=offwhite}
\setbeamercolor{alerted text}{fg=offwhite}
\setbeamercolor{example text}{fg=offwhite}
\setbeamercolor{section in toc}{fg=offwhite}
\setbeamercolor{subsection in toc}{fg=offwhite}
\setbeamercolor{subsubsection in toc}{fg=offwhite}
\setbeamercolor{caption}{fg=offwhite}
\setbeamercolor{caption name}{fg=offwhite}
\setbeamercolor{footnote mark}{fg=offwhite}
\setbeamercolor{block title}{fg=offwhite,bg=softblack}
\setbeamercolor{block body}{fg=offwhite,bg=softblack}
\setbeamercolor{block title alerted}{fg=offwhite,bg=softblack}
\setbeamercolor{block body alerted}{fg=offwhite,bg=softblack}
\setbeamercolor{block title example}{fg=offwhite,bg=softblack}
\setbeamercolor{block body example}{fg=offwhite,bg=softblack}
\setbeamercolor{bibliography entry author}{fg=offwhite}
\setbeamercolor{bibliography entry title}{fg=offwhite}
\setbeamercolor{bibliography entry location}{fg=offwhite}
\setbeamercolor{bibliography entry note}{fg=offwhite}
\hypersetup{colorlinks=true,urlcolor=offwhite,citecolor=offwhite,linkcolor=offwhite,filecolor=offwhite,pdfdisplaydoctitle=true}

% Reusable text and table commands.
\newcommand{\samelineand}{\qquad}
\newcolumntype{Y}{>{\centering\arraybackslash}X}

% \begin{SlideContents}...\end{SlideContents}
\newenvironment{SlideContents}{%
    \tableofcontents[subsectionstyle=show]
}{}

% \begin{SlideBulletsPhoto}{file}{height}{caption}...bullets...\end{SlideBulletsPhoto}
\newenvironment{SlideBulletsPhoto}[3]{%
    \def\slidebulletsphotofile{#1}%
    \def\slidebulletsphotoheight{#2}%
    \def\slidebulletsphotocaption{#3}%
    \begin{columns}[c,onlytextwidth]
    \begin{column}{0.40\textwidth}
}{%
    \end{column}%
    \begin{column}{0.54\textwidth}
        \begin{figure}
            \centering
            \presentationimage{\slidebulletsphotofile}{\slidebulletsphotoheight}{\linewidth}
            \caption{\slidebulletsphotocaption}
        \end{figure}
    \end{column}%
    \end{columns}%
}

% \begin{SlideTitlePhoto}{file}{height}{maximum width}...caption...\end{SlideTitlePhoto}
\newenvironment{SlideTitlePhoto}[3]{%
    \begin{figure}
        \centering
        \presentationimage{#1}{#2}{#3}%
}{\end{figure}}

% \begin{SlidePhotoPair}...\SlidePhotoPanel{file}{height}{caption}...\end{SlidePhotoPair}
\newenvironment{SlidePhotoPair}{%
    \begin{figure}\centering
}{\end{figure}}
\newcommand{\SlidePhotoPanel}[3]{%
    \begin{subfigure}[t]{0.46\linewidth}
        \vspace{0pt}
        \centering
        \presentationimage{#1}{#2}{\linewidth}
        \caption{#3}
    \end{subfigure}%
}

% Use native column environments inside SlideThreeColumns for each fixed-width column.
\newenvironment{SlideThreeColumns}{%
    \begin{columns}[T,onlytextwidth]
}{\end{columns}}

% \begin{SlideTable}...caption and tabular...\end{SlideTable}
\newenvironment{SlideTable}{%
    \begin{table}\centering
}{\end{table}}

% Use two native column environments inside SlideCodeExplanation.
\newenvironment{SlideCodeExplanation}{%
    \begin{columns}[c,onlytextwidth]
}{\end{columns}}

% Keep this command active through the end of a timeline frame.
\newcommand{\SlideTimeline}{\centering}

% \begin{SlideReferences}...bibliography...\end{SlideReferences}
\newenvironment{SlideReferences}{%
    \footnotesize\renewcommand{\bibsection}{}
}{}

% \primaryauthor{name} identifies and underlines the main author.
\newcommand{\primaryauthor}[1]{\underline{#1}}

% \presentationimage{file}{height}{maximum width}
\newsavebox{\presentationimagebox}
\newlength{\presentationimagemaxwidth}
\newcommand{\presentationimage}[3]{%
    \begingroup
    \sbox{\presentationimagebox}{\includegraphics[height=#2]{#1}}%
    \setlength{\presentationimagemaxwidth}{#3}%
    \ifdim\wd\presentationimagebox>\presentationimagemaxwidth
        \begin{tikzpicture}[baseline=0pt]
            \path[use as bounding box] (-.5\presentationimagemaxwidth,0)
                rectangle (.5\presentationimagemaxwidth,#2);
            \clip (-.5\presentationimagemaxwidth,0) rectangle (.5\presentationimagemaxwidth,#2);
            \node[anchor=south,inner sep=0pt,outer sep=0pt] at (0,0) {\usebox{\presentationimagebox}};
        \end{tikzpicture}%
    \else
        \usebox{\presentationimagebox}%
    \fi
    \endgroup
}

% List markers and internal Beamer adjustments.
\setbeamertemplate{itemize item}[circle]
\setbeamertemplate{itemize subitem}[circle]
\setbeamertemplate{itemize subsubitem}[circle]
\makeatletter
% Keep author-year citations linked across line breaks.
\def\hyper@natlinkbreak#1#2{\hyper@linkend#1\hyper@linkstart{cite}{cite.#2}}
\xpatchcmd{\itemize}
    {\def\makelabel}
    {\ifnum\@itemdepth=1\relax\setlength\itemsep{1.6ex}\else\setlength\itemsep{0.4ex}\fi\def\makelabel}
    {}{\PackageError{presentation}{Could not patch list spacing}{Check the installed Beamer version.}}
% Clear the previous title so untitled slides can use section headings.
\renewcommand*\beamer@checkframetitle{
    \global\let\beamer@frametitle\relax
    \@ifnextchar\bgroup\beamer@inlineframetitle{}
}
\AtBeginPart{\beamer@tocsectionnumber=0\relax\setcounter{section}{0}}
% Record the next slide number in each table-of-contents entry.
\newcounter{tocframenumber}
\newcommand{\tocdots}{\leavevmode\cleaders\hb@xt@ 0.8em{\hss.\hss}\hfill\kern\z@}
\patchcmd{\beamer@section}
    {\addtocontents{toc}{\protect\beamer@sectionintoc{\the\c@section}{#2}{\the\c@page}{\the\c@part}{\the\beamer@tocsectionnumber}}}
    {\setcounter{tocframenumber}{\value{framenumber}}\stepcounter{tocframenumber}\addtocontents{toc}{\protect\beamer@sectionintoc{\the\c@section}{#2\tocdots\thetocframenumber}{\the\c@page}{\the\c@part}{\the\beamer@tocsectionnumber}}}
    {}{\PackageError{presentation}{Could not patch section entries}{Check the installed Beamer version.}}
\patchcmd{\beamer@subsection}
    {\addtocontents{toc}{\protect\beamer@subsectionintoc{\the\c@section}{\the\c@subsection}{#2}{\the\c@page}{\the\c@part}{\the\beamer@tocsectionnumber}}}
    {\setcounter{tocframenumber}{\value{framenumber}}\stepcounter{tocframenumber}\addtocontents{toc}{\protect\beamer@subsectionintoc{\the\c@section}{\the\c@subsection}{#2\tocdots\thetocframenumber}{\the\c@page}{\the\c@part}{\the\beamer@tocsectionnumber}}}
    {}{\PackageError{presentation}{Could not patch subsection entries}{Check the installed Beamer version.}}
\patchcmd{\beamer@subsubsection}
    {\addtocontents{toc}{\protect\beamer@subsubsectionintoc{\the\c@section}{\the\c@subsection}{\the\c@subsubsection}{#2}{\the\c@page}{\the\c@part}{\the\beamer@tocsectionnumber}}}
    {\setcounter{tocframenumber}{\value{framenumber}}\stepcounter{tocframenumber}\addtocontents{toc}{\protect\beamer@subsubsectionintoc{\the\c@section}{\the\c@subsection}{\the\c@subsubsection}{#2\tocdots\thetocframenumber}{\the\c@page}{\the\c@part}{\the\beamer@tocsectionnumber}}}
    {}{\PackageError{presentation}{Could not patch subsubsection entries}{Check the installed Beamer version.}}
\ifnotes
    % The slide already has page anchors; do not duplicate them on the notes page.
    \xpatchcmd{\beamer@outsideframenote}
        {\begingroup}
        {\begingroup\let\Hy@EveryPageAnchor\@empty}
        {}{\PackageWarning{presentation}{Could not suppress duplicate notes anchors}}
\fi
\makeatother

% Captions, navigation, and slide layout.
\setlength\belowdisplayshortskip{1ex}
\setlength{\abovecaptionskip}{1ex}
\setlength{\belowcaptionskip}{0.5ex}
\setbeamertemplate{caption}[numbered]
\setbeamertemplate{caption label separator}[period]
\setbeamertemplate{navigation symbols}{}
\setbeamertemplate{blocks}[default]
\setbeamertemplate{section in toc}[sections numbered]
\setbeamertemplate{subsection in toc}{
    \leavevmode\leftskip=3em
    \rlap{\hskip-1.8em\inserttocsectionnumber.\inserttocsubsectionnumber}
    \inserttocsubsection\par\vskip0.3em
}
\setbeamertemplate{subsubsection in toc}{
    \leavevmode\leftskip=5em
    \rlap{\hskip-2.8em\inserttocsectionnumber.\inserttocsubsectionnumber.\inserttocsubsubsectionnumber}
    \inserttocsubsubsection\par
}

\newsavebox{\presentationtitlebox}
\setbeamertemplate{title page}{
    \vfill
    \begin{flushleft}
        {\usebeamerfont{title}\usebeamercolor[fg]{title}%
            \sbox{\presentationtitlebox}{\scalebox{1.08}{\inserttitle}}%
            \ifdim\wd\presentationtitlebox>\linewidth
                \sbox{\presentationtitlebox}{\parbox{\linewidth}{\raggedright\inserttitle}}%
                \ifdim\dimexpr\ht\presentationtitlebox+\dp\presentationtitlebox\relax>.4\paperheight
                    \PackageWarning{presentation}{Long title reduced to fit the page}%
                    \resizebox*{!}{.4\paperheight}{\usebox{\presentationtitlebox}}\par
                \else
                    \usebox{\presentationtitlebox}\par
                \fi
            \else
                \usebox{\presentationtitlebox}\par
            \fi}
        \medskip
        {\usebeamerfont{subtitle}\insertsubtitle\par}
        \bigskip
        {\usebeamerfont{author}\insertauthor\par}
        \medskip
        {\usebeamerfont{institute}\insertinstitute\par}
        \medskip
        {\usebeamerfont{date}\insertdate\par}
    \end{flushleft}
    \vfill
}
\setbeamertemplate{frametitle}{
    \nointerlineskip\vskip4ex
    {\usebeamerfont{frametitle}\usebeamercolor[fg]{frametitle}\insertframetitle\par}
    \ifx\insertframesubtitle\empty\else
        {\usebeamerfont{framesubtitle}\usebeamercolor[fg]{framesubtitle}\insertframesubtitle\par}
    \fi
    \vskip1.5ex
}
\addtobeamertemplate{frametitle}{
    \ifx\insertframetitle\empty
        \frametitle{\insertsectionhead}
        \ifx\insertframesubtitle\empty
            \framesubtitle{\insertsubsectionhead}
        \fi
    \fi
}{}
\setbeamertemplate{footline}{
    \usebeamerfont{footline}\usebeamercolor[fg]{page number in head/foot}
    \hbox to \paperwidth{\hskip14mm\strut\hfill\insertframenumber{} / \inserttotalframenumber\hskip14mm}
    \vskip2ex
}

% The notes entry point enables a second screen on the right.
\ifnotes
    \setbeamerfont{note page}{size=\Large}
    \setbeamercolor{note page}{fg=offwhite}
    \newsavebox{\presentationnotebox}
    \setbeamertemplate{note page}{%
        \sbox{\presentationnotebox}{%
            \begin{tikzpicture}
                \node[align=center,text width=\dimexpr\paperwidth-28mm\relax]
                    {\hyphenpenalty=10000\exhyphenpenalty=10000\insertnote\par};
            \end{tikzpicture}%
        }%
        \begin{tikzpicture}[remember picture,overlay]
            \fill[softblack] (current page.south west) rectangle (current page.north east);
            \node[inner sep=0pt] at (current page.center) {%
                \ifdim\dimexpr\ht\presentationnotebox+\dp\presentationnotebox\relax>\dimexpr\paperheight-20mm\relax
                    \PackageWarning{presentation}{Long speaker note reduced to fit the page}%
                    \resizebox*{!}{\dimexpr\paperheight-20mm\relax}{\usebox{\presentationnotebox}}%
                \else
                    \usebox{\presentationnotebox}%
                \fi
            };
        \end{tikzpicture}%
    }
    \setbeameroption{show notes on second screen=right}
\fi

\pdfstringdefDisableCommands{\def\inst#1{}\def\and{, }\def\underline#1{#1}\def\primaryauthor#1{#1}\def\samelineand{; }}

% Localized section labels follow the explicitly selected document language.
% Imported Babel languages can share a hyphenation identifier, so compare names.
\newcommand{\ifpresentationlanguage}[3]{%
    \ifnum\pdfstrcmp{\languagename}{#1}=0 #2\else #3\fi
}
\newcommand{\contentslabel}{%
    \ifpresentationlanguage{brazilian}{Sumário}{%
        \ifpresentationlanguage{spanish}{Índice}{Contents}%
    }%
}
\newcommand{\referenceslabel}{%
    \ifpresentationlanguage{brazilian}{Referências}{%
        \ifpresentationlanguage{spanish}{Referencias}{References}%
    }%
}

% Format the presentation date in the selected document language.
\newcommand{\presentationmonthname}{%
    \ifpresentationlanguage{english}{%
        \ifcase\presentationdatemonth\or January\or February\or March\or April\or May\or June\or July\or August\or September\or October\or November\or December\fi
    }{\ifpresentationlanguage{brazilian}{%
        \ifcase\presentationdatemonth\or janeiro\or fevereiro\or março\or abril\or maio\or junho\or julho\or agosto\or setembro\or outubro\or novembro\or dezembro\fi
    }{%
        \ifcase\presentationdatemonth\or enero\or febrero\or marzo\or abril\or mayo\or junio\or julio\or agosto\or septiembre\or octubre\or noviembre\or diciembre\fi
    }}%
}
\newcommand{\presentationdate}{%
    \ifpresentationlanguage{english}{%
        \presentationmonthname\space\presentationdateday, \presentationdateyear
    }{%
        \presentationdateday\space de\space\presentationmonthname\space de\space\presentationdateyear
    }%
}

